\documentclass[11pt]{article}
\usepackage[margin=1in]{geometry}
\usepackage{amsmath,amssymb,amsthm}
\usepackage{hyperref}

\newtheorem{theorem}{Theorem}
\newtheorem{lemma}[theorem]{Lemma}
\theoremstyle{remark}
\newtheorem*{remark}{Remark}

\title{Commutativity of Submonoid Quantales of Ring and Semiring Domains}
\author{Salem Levinstein}
\date{November 26, 2026}

\begin{document}
\maketitle

\begin{abstract}
We resolve the first part of Question~34 of Mesyan (\texttt{arXiv:2611.31109v1}), which asks whether, for a domain $R$, commutativity of $R$ is equivalent to commutativity of its quantale $\mathcal{M}(R)$ of additive submonoids. The answer splits between rings and semirings: for every (possibly nonunital) ring without zero divisors, $\mathcal{M}(R)$ is commutative if and only if $R$ is commutative; for semiring domains, the equivalence fails, as shown by a $14$-element noncommutative semiring domain $R_{14}$ whose submonoid quantale $\mathcal{M}(R_{14})$ is commutative. All results are formally verified in Lean~4.
\end{abstract}

\section{Introduction}
In~\cite{Mesyan2026}, Mesyan studies the quantale $\mathcal{M}(R)$ of all additive submonoids of a (possibly nonunital) semiring $R$, equipped with the multiplication
\[
  S \cdot T = \left\{\sum_{i=1}^{k} s_i t_i : k \ge 0,\; s_i \in S,\; t_i \in T\right\}
\]
(Definition~4 of~\cite{Mesyan2026}). Here a \emph{semiring} consists of a commutative monoid $(R,+,0)$ and a semigroup $(R,\cdot)$ distributing over $+$ on both sides with $0r = r0 = 0$; a \emph{domain} is a nonzero semiring such that $rp \ne 0$ for all nonzero $r,p \in R$ (Remark~14). Section~9 of~\cite{Mesyan2026} poses the following open problem:

\medskip
\noindent\textbf{Question 34 (\cite{Mesyan2026}).} \emph{For a domain $R$, is $R$ being commutative equivalent to $\mathcal{M}(R)$ being commutative? (See Proposition~22~(3) and subsequent comments.) More broadly, for which (semi)rings $R$ is $\mathcal{M}(R)$ commutative?}
\medskip

Proposition~22(3) of~\cite{Mesyan2026} established the equivalence under the additive hypothesis that every $r \in R$ satisfies $r = 3r$ or $r \ne nr$ for all $n > 1$, leaving open prime-characteristic ring domains with $p \ge 3$ as well as general semiring domains. We settle the first sentence of Question~34 completely.

\section{Reduction to Cyclic Submonoids}
For $x \in R$, let $\langle x \rangle = \{nx : n \in \mathbb{N}\}$ denote the cyclic additive submonoid generated by $x$.

\begin{lemma}[Reduction (R)]\label{lem:R}
For any semiring $R$, $\mathcal{M}(R)$ is commutative if and only if $\langle rs \rangle = \langle sr \rangle$ for all $r, s \in R$.
\end{lemma}

\begin{proof}
For all $r, s \in R$, $\langle r \rangle \cdot \langle s \rangle = \langle rs \rangle$ (Lemma~21(1) of~\cite{Mesyan2026}), because $(nr)(ms) = (nm)(rs) \in \langle rs \rangle$ while $rs \in \langle r \rangle \cdot \langle s \rangle$. Thus commutativity of $\mathcal{M}(R)$ implies $\langle rs \rangle = \langle sr \rangle$. Conversely, if $\langle rs \rangle = \langle sr \rangle$ for all $r, s \in R$, then for any $S, T \in \mathcal{M}(R)$ and any $s \in S$, $t \in T$, we have $st \in \langle st \rangle = \langle ts \rangle \subseteq T \cdot S$. Hence $S \cdot T \subseteq T \cdot S$, and symmetry gives $S \cdot T = T \cdot S$.
\end{proof}

\section{Affirmative Answer for Ring Domains}

\begin{theorem}\label{thm:A}
Let $R$ be a (possibly nonunital) ring without zero divisors. Then $\mathcal{M}(R)$ is commutative if and only if $R$ is commutative.
\end{theorem}

\begin{proof}
If $R$ is commutative, then $\mathcal{M}(R)$ is commutative by Lemma~\ref{lem:R}. Conversely, assume $\mathcal{M}(R)$ is commutative and suppose for contradiction that $ab \ne ba$ for some $a, b \in R$. Then $a \ne 0$ and $b \ne 0$, so $ab \ne 0$ and $ba \ne 0$.

By Lemma~\ref{lem:R}, $\langle ab \rangle = \langle ba \rangle$, so there exist integers $n, m \ge 1$ with $ba = n(ab)$ and $ab = m(ba)$. Then $(mn - 1)(ab) = 0$. Since $ab \ne ba$, we have $n \ne 1$, hence $mn - 1 \ge 1$. Let $p \ge 1$ be the least positive integer with $p(ab) = 0$. For every $y \in R$, $(ab)(py) = (p(ab))y = 0$, whence $py = 0$ because $ab \ne 0$ and $R$ has no zero divisors. If $p = uv$ with $u, v \ge 1$, then $(u(ab))(v(ab)) = p((ab)^2) = 0$, so $u(ab) = 0$ or $v(ab) = 0$; by minimality of $p$ and $ab \ne 0$, $p$ is prime, and for any $y \ne 0$ and $z \in \mathbb{Z}$, $zy = 0$ if and only if $p \mid z$.

Let $\lambda = n \bmod p \in \mathbb{F}_p$. Since $ba \ne 0$ and $ba \ne ab$, we have $\lambda \notin \{0, 1\}$. Moreover,
\[
  ba^2 = (ba)a = \lambda(ab)a = \lambda a(ba) = \lambda^2 a^2 b.
\]
Applying Lemma~\ref{lem:R} to $(a + a^2, b)$ yields $\mu \in \mathbb{F}_p$ with $b(a + a^2) = \mu(a + a^2)b$, which expands to
\[
  \bigl((\lambda - \mu)a + (\lambda^2 - \mu)a^2\bigr)b = (\lambda - \mu)ab + (\lambda^2 - \mu)a^2 b = 0.
\]
Cancelling $b \ne 0$ gives $(\lambda - \mu)a + (\lambda^2 - \mu)a^2 = 0$. If $\lambda^2 - \mu = 0$ in $\mathbb{F}_p$, then $(\lambda - \mu)a = 0$, so $\lambda = \mu = \lambda^2$, contradicting $\lambda \notin \{0, 1\}$. Therefore $\lambda^2 - \mu \ne 0$, and $a^2 = ca$ for $c = (\mu - \lambda)(\lambda^2 - \mu)^{-1} \in \mathbb{F}_p$. Since $a^2 \ne 0$, we have $c \ne 0$.

Set $e = c^{-1}a \in R$. Then $e \ne 0$ and $e^2 = c^{-2}a^2 = c^{-1}a = e$. For every $y \in R$,
\[
  e(ey - y) = e^2 y - ey = 0 \implies ey = y, \qquad (ye - y)e = ye^2 - ye = 0 \implies ye = y.
\]
Thus $e$ is a two-sided identity of $R$. Consequently $a = ce$ is central:
\[
  ab = (ce)b = c(eb) = cb = c(be) = b(ce) = ba,
\]
contradicting $ab \ne ba$.
\end{proof}

\begin{remark}
Example~6 of~\cite{Mesyan2026} ($X(\mathbb{F}_3)$) is a noncommutative ring with commutative $\mathcal{M}(X(\mathbb{F}_3))$ and zero divisors, showing that the no-zero-divisors hypothesis in Theorem~\ref{thm:A} is essential.
\end{remark}

\section{Counterexample for Semiring Domains}

\begin{theorem}\label{thm:B}
There exists a $14$-element noncommutative semiring domain $R_{14}$ such that $\mathcal{M}(R_{14})$ is commutative.
\end{theorem}

\begin{proof}
Let $G = \mathbb{Z}_2^2$ and set $R_{14} = \{0, \infty\} \cup (\mathbb{Z}_3 \times G)$, so $|R_{14}| = 14$. Define the asymmetric bicharacter $\beta : G \times G \to \mathbb{Z}_3^\times$ by $\beta((g_1, g_2), (h_1, h_2)) = (-1)^{g_1 h_2}$. Define addition on $R_{14}$ by taking $0$ as the additive identity, $\infty + x = x + \infty = \infty$ for all $x \ne 0$, and
\[
  (i, g) + (j, h) = \begin{cases} (i + j \bmod 3,\, g) & \text{if } g = h, \\ \infty & \text{if } g \ne h. \end{cases}
\]
Define multiplication on $R_{14}$ by taking $0$ to be absorbing, $\infty \cdot x = x \cdot \infty = \infty$ for all $x \ne 0$, and
\[
  (i, g) \cdot (j, h) = \bigl(ij\beta(g, h) \bmod 3,\; g + h\bigr).
\]
Addition is commutative and associative with identity $0$, since any sum involving $\infty$ or two distinct fibres collapses to $\infty$. Multiplication is associative because $\beta$ is biadditive and hence satisfies the $2$-cocycle identity $\beta(g,h)\beta(g+h,k) = \beta(g,h+k)\beta(h,k)$. Two-sided distributivity follows from $\mathbb{Z}_3$-bilinearity inside fibres together with the fact that translation $h \mapsto g + h$ is a bijection on $G$.

No product of nonzero elements equals $0$ (note that $(0,g) \in \mathbb{Z}_3 \times G$ is distinct from the additive identity $0$), so $R_{14}$ is a domain. It is noncommutative because
\[
  (1, (1,0)) \cdot (1, (0,1)) = (2, (1,1)) \ne (1, (1,1)) = (1, (0,1)) \cdot (1, (1,0)).
\]
Finally, for any $r, s \in R_{14}$, if $r$ or $s$ lies in $\{0,\infty\}$ then $rs = sr$. If $r = (i, g)$ and $s = (j, h)$, then $rs$ and $sr$ lie in the same fibre $g + h$ with first coordinates differing by a factor in $\mathbb{Z}_3^\times = \{1, 2\}$, so either $rs = sr$ or $rs = 2(sr)$ and $sr = 2(rs)$. Thus $\langle rs \rangle = \langle sr \rangle$ for all $r, s \in R_{14}$, and Lemma~\ref{lem:R} implies that $\mathcal{M}(R_{14})$ is commutative.
\end{proof}

\begin{remark}
The construction of Theorem~\ref{thm:B} generalizes to any modulus $d \ge 3$ and any non-symmetric bicharacter $G \times G \to \mathbb{Z}_d^\times$. Every fibre element $r = (1, g) \in R_{14}$ satisfies $4r = r \ne 3r$, lying outside the scope of Proposition~22(3) of~\cite{Mesyan2026}.
\end{remark}

\section{Formal Verification}
Lemma~\ref{lem:R}, Theorem~\ref{thm:A}, and Theorem~\ref{thm:B} have been formally verified in Lean~4 with Mathlib (file \texttt{Mesyan\-Submonoid\-Quantale.lean}), using only the standard axioms \texttt{propext}, \texttt{Classical.choice}, and \texttt{Quot.sound}, with zero \texttt{sorry} and zero warnings. A standalone Python script (\texttt{verify.py}) independently verifies the semiring and domain axioms of $R_{14}$ and checks commutativity of $\mathcal{M}(R_{14})$ over all $90$ additive submonoids.

\begin{thebibliography}{1}
\bibitem{Mesyan2026}
Z.~Mesyan, \emph{Quantales arising from rings}, arXiv:2611.31109v1 [math.RA], November 23, 2026.
\end{thebibliography}

\end{document}
